\chapter*{Summary}
\addcontentsline{toc}{chapter}{Summary}
\markboth{Summary}{Summary}

Around 2000, Background Oriented Schlieren proposed to replace the large optical elements used by traditional schlieren with a background pattern, digital optical sensors and the use of cross-correlation algorithms \cite{Dalziel2000}. The cost was an inherent tie between the sensitivity and the physical distances in the optical setup, which lead to a design compromise between sensitivity and resolution \cite{Raffel2015, Schmidt2025}. By replacing the printed background pattern with a laser-speckled pattern that constraint is relaxed \cite{Meier&Roesgen2013}. Since the speckle is always in-focus, it can be positioned anywhere in the line of sight of the camera, allowing for unusual focus positions, namely in front of or inside of it. The resulting technique, Laser-Speckled Background Oriented Schlieren, remains relatively unexplored compared to standard BOS and lacks a coherent design framework for its implementation. This thesis investigates the application of the method to a supersonic wind tunnel, through four experimental campaigns that address each of the three components of the method: the optical configuration, the speckle generator and the refractive object. The last of these campaigns brings the three together in a high-speed wind tunnel.

A design methodology was first developed to solve the optical setup performance parameters for a set of physical variables and correctly size the experimental setups. That lead to the identification of the variables that influence the sensitivity-resolution ratio and allows for improved performance of SBOS compared to BOS. It also gave rise to surface maps that permitted the main trends of the design space to be identified and transformed into coherent design guidelines using a spatial resolution criterion. An angled glass slide that induced a known displacement served to validate the methodology and compare different in-focus positions, reproducing the predicted sensitivities within 5\% at all but one point. Uncertainty budgets were built for the predicted and the measured displacements. A dedicated campaign on the speckle pattern established that the f-stop can be used to control the average speckle size while the optical magnification cannot, which contradicts the Rayleigh limit of Equation \ref{equation: Rayleigh Limit}. Nevertheless, it led to a generator capable of producing speckle sizes between 2 and 5 pixels. An under-expanded compressible jet was then imaged and its displacement field converted into a quantitative density field through Abel inversion with Tikhonov regularization \cite{Daun2006}. The results agree with the shock cell correlations of Emden \cite{Emden1899} and of Hartmann and Lazarus \cite{Hartmann1941} and, over the first shock cell, with the centerline density of Panda and Seasholtz \cite{Panda1999}.

The final experimental campaign implemented the method to image a linear aerospike nozzle flow in TU Delft's ST-15 supersonic wind tunnel. With the focus inside the test section the usual model $S = l_{eff} M_{ag}$ no longer holds, so a method was introduced that recovers the sensitivity from the ratio of displacements measured at two configurations, which agreed with the geometric sensitivity within 5\% wherever both could be computed. In-focus SBOS presented similar spatial resolution to schlieren images of the same flow, capable of identifying the same compression waves, expansion fans and shear layer. Whereas the defocused configurations blurred the pattern to decorrelation under strong density gradients. Taken together, the four campaigns propose a design route for SBOS setups and show that the method can be applied in a high-speed wind tunnel with a simple optical arrangement and a low power laser.